% from ega4-21c-fr.tex
\oldpage[IV-4]{311}
(iii) Sans supposer que $Y$ soit localement factoriel, mais en supposant par contre que
$X$ est intersection complète relativement à $Y$ (chap.~V), nous verrons au chap.~V
qu'on a un résultat plus précis que (21.12.12), en exprimant $T$ comme cycle
1-codimensionnel d'une section d'un $\mathcal O_X$-Module inversible, canoniquement
associée à $g$. Ceci s'appliquera notamment quand $X$ et $Y$ sont tous deux réguliers,
ou lorsque $g$ est un $S$-morphisme, $X$ et $Y$ étant des préschémas lisses sur $S$.

(iv) On peut se demander si (21.12.10) admet une réciproque; autrement dit, pour un
anneau local noethérien donné $A$, si l'on pose $Y=\operatorname{Spec}(A)$, $y$
désignant le point fermé de $Y$, et si, pour tout préschéma localement noethérien $X$ et
tout morphisme $g:X\to Y$ tel que $g':X'\to Y'$ soit une immersion ouverte, toute
composante irréductible de $X_y$ est de codimension $\leq1$ dans $X$ ou est réduite à
un point, alors est-il vrai que $A$ vérifie la condition (W) de (21.12.8)?

(v) Soient $Y$ un préschéma intègre \emph{régulier}, $X$ un préschéma localement
noethérien \emph{normal}, $g:X\to Y$ un morphisme localement de type fini; supposons
en outre que, si $\eta$ est le point générique de $Y$, la fibre $X_\eta$ soit \emph{étale}
sur $k(\eta)$, et soit $T'$ le complémentaire du plus grand ouvert $V'$ de $X$ tel que
$g|V':V'\to Y$ soit étale; est-il vrai alors que toutes les composantes irréductibles de
$T'$ soient de codimension $1$? Il en est bien ainsi lorsqu'on suppose en outre que $g$
est \emph{localement quasi-fini} (« théorème de pureté » de Zariski--Nagata). On peut
montrer que la conjecture précédente serait conséquence de la suivante (et lui serait
même équivalente si la conjecture de (iv) était vérifiée): pour un anneau local régulier
$A$ contenu dans un anneau local intègre et intégralement clos $B$, tel que $B$ soit un
$A$-module fini, l'ouvert $U$ des points de $\operatorname{Spec}(B)$ en lesquels
$\operatorname{Spec}(B)$ est non ramifié sur $\operatorname{Spec}(A)$ (ou, ce qui
revient au même par (18.10.1), étale sur $\operatorname{Spec}(A)$) est affine.
\end{remarks}

\begin{proposition}[21.12.15]
\label{IV.21.12.15}
Soient $S$ un préschéma, $g:X\to S$ un morphisme plat et localement de présentation
finie, dont les fibres $X_s=g^{-1}(s)$ sont géométriquement irréductibles (4.5.2),
$h:Y\to S$ un morphisme lisse; on note $Y_s=h^{-1}(s)$ les fibres de ce morphisme.
Soit $f:X\to Y$ un $S$-morphisme propre tel que, pour tout point maximal $\eta$ de
$S$, le morphisme $f_\eta:X_\eta\to Y_\eta$ soit un isomorphisme. Alors $f$ est un
isomorphisme.
\end{proposition}

\begin{proof}
Puisque $h$ est plat, tout point maximal de $Y$ est au-dessus d'un point maximal de $S$
(2.3.4), donc la réunion des $Y_\eta$, lorsque $\eta$ parcourt l'ensemble des points
maximaux de $S$, est dense dans $Y$; comme les $f_\eta$ sont des isomorphismes, $f$
est dominant et par suite \emph{surjectif} puisqu'il est propre. Il suffit donc de montrer
que $f$ est une \emph{immersion ouverte}. Compte tenu de (17.9.5), il suffit de prouver
que pour tout $s\in S$, $f_s:X_s\to Y_s$ est une immersion ouverte.

Comme tout $s\in S$ est générisation d'un point maximal $\eta$, on peut déjà, en
faisant le changement de base $\operatorname{Spec}(\mathcal O_{S',s})\to S$, où $S'$
est le sous-préschéma réduit de $S$ ayant $\overline{\{\eta\}}$ comme espace
sous-jacent, se ramener au cas où $S$ est un schéma local et intègre de point fermé $s$:
les fibres aux points $\eta$ et $s$ ne sont en effet pas changées et les propriétés des
morphismes $g$, $h$ et $f$ se conservent par changement de base. En outre, la question
est locale sur $Y$; remplaçant $Y$ par un voisinage ouvert affine $U$ dans $Y$ d'un point
de $Y_s$, et $X$ par $f^{-1}(U)$, qui est quasi-compact puisque $f$ est propre, on voit
qu'on peut supposer en outre que $X$ et $Y$ sont de présentation finie sur $S$.

Il existe alors un sous-anneau local noethérien $A_0$ de
$A=\mathcal O_{S,s}$ tel que $A_0\to A$ soit un homomorphisme local, deux morphismes
de type fini $g_0:X_0\to S_0=\operatorname{Spec}(A_0)$,
$h_0:Y_0\to S_0$ et un $S_0$-morphisme
